% =============================================================================
% 2023-2026
% =============================================================================

\section{Antagonist}

% Bu bolumde birden fazla olay var: 2023 genel secimi, lider degisimi,
% 2024 yerel secimi, Bahceli cagrisi, Imamoglu tutuklamasi, Istanbul kayyumu.
%
% Kritik ayrim, ikisinin AYNI ANDA olmadigi:
%   Yolsuzluk ve cete   kirilma 2023-11  (lider degisimiyle ayni ay)
%   Erdogan ve AKP      kirilma 2024-09  (on ay sonra)
% Saf uslup etkisi olsaydi ikisi de 2023-11'de kirilirdi.
%
% Son uc olayda korpus 24 ayi doldurmuyor: donus GOZLEM DISI, "gerceklesmedi"
% degil.

\subsection*{Erdoğan ve AKP}

2023 Mayıs genel seçimlerinden sonraki dönemde Erdoğan ve AKP kelimelerinin frekansları
sürekli olarak artışa geçiyor ve İmamoğlu tutuklanmasından sonraki üç aylık periyotta pik yapıyor.
Buna karşı olarak darbe girişimi ve referandum sonrasında aynı tema sönümleniyor.
Düşüş oranları benzer olsa da eski seviyeye dönüş darbe girişiminde 16, referandumda 5 ay sürüyor.

\begin{table}[H]
\centering
\caption{Erdoğan ve AKP: olay çevresindeki hareket}
\label{tab:olay-antagonist_akponly}
\small
\begin{tabular}{@{}l r r r c l@{}}
\toprule
Olay & Taban & Sonrası & Değişim & & Dönüş \\
\midrule
  Darbe girişimi, Tem. 2016 & 16.5 & 9.7 & -41\% & \azalan & 16 ay \\
  Referandum, Nis. 2017 & 11.2 & 7.1 & -37\% & \azalan & 5 ay \\
  Genel seçim, Haz. 2018 & 12.6 & 14.9 & +18\% & \artan & --\\
  Yerel seçim, Mar. 2019 & 15.2 & 24.1 & +58\% & \artan & -- \\
  Genel seçim, May. 2023 & 9.9 & 18.8 & +90\% & \artan & -- \\
  Yerel seçim, Mar. 2024 & 21.9 & 31.8 & +45\% & \artan & -- \\
  Bahçeli çağrısı, Eki. 2024 & 31.4 & 43.0 & +37\% & \artan & -- \\
  İmamoğlu tutuklaması, Mar. 2025 & 36.3 & 53.1 & +46\% & \artan & -- \\
  İstanbul kayyumu, Eyl. 2025 & 49.7 & 50.0 & +1\% & \duraga & -- \\
\bottomrule
\end{tabular}
\kaynaknot{Taban olaydan önceki 12 ayın, \emph{Sonrası} sonraki üç ayın
ortalamasıdır; ikisi de 10.000 kelimede vukuat sayısıdır. \emph{Dönüş},
kaydın tabanın \%90'ına döndüğü ay sayısıdır. Korpus 31 Mart 2026'da bittiği
için son üç olayda 24 aylık pencere dolmuyor: oralarda dönüş gözlem dışıdır,
gerçekleşmediği anlamına gelmez.}
\end{table}


\subsection*{Yolsuzluk ve Çete}

Kılıçdaroğlu döneminde anayasa referandumundan sonra tırmanmaya başlayan yolsuzluk ve çete teması 
İBB seçimlerinin yenilendiği 2019 yerel seçimleri sonrasındaki 3 ay pencere içerisinde $-$46\% oranında azalıyor.
13 ay sonra eski seviyesine dönse de çete, mafya, baron ve rüşvet gibi kelimeler lider değişimiyle birlikte 
söylemden büyük oranda çekiliyor ($-$66\%). [ÖZEL DÖNEMİNDEKİ İKİ ARTIŞ MUHTEMELEN ADALET BAKANLIĞINDAN
KAYNAKLANIYOR. AMA GÖRMEDEN GEÇEMEM. VERİ İLE KONTROL EDİP BURAYA DÖN]




\begin{table}[H]
\centering
\caption{Yolsuzluk ve çete: olay çevresindeki hareket}
\label{tab:olay-antagonist_general}
\small
\begin{tabular}{@{}l r r r c l@{}}
\toprule
Olay & Taban & Sonrası & Değişim & & Dönüş \\
\midrule
  Darbe girişimi, Tem. 2016 & 3.9 & 2.5 & -34\% & \azalan & 10 ay \\
  Referandum, Nis. 2017 & 2.0 & 4.5 & +128\% & \artan & -- \\
  Genel seçim, Haz. 2018 & 11.8 & 29.6 & +151\% & \artan & -- \\
  Yerel seçim, Mar. 2019 & 21.4 & 11.6 & -46\% & \azalan & 13 ay \\
  Genel seçim, May. 2023 & 25.0 & 22.5 & -10\% & \duraga & -- \\
  Yerel seçim, Mar. 2024 & 14.9 & 4.4 & -71\% & \azalan & >24 ay \\
  Bahçeli çağrısı, Eki. 2024 & 4.7 & 3.4 & -27\% & \azalan & 4 ay \\
  İmamoğlu tutuklaması, Mar. 2025 & 3.8 & 5.5 & +47\% & \artan & -- \\
  İstanbul kayyumu, Eyl. 2025 & 5.7 & 7.8 & +36\% & \artan & -- \\
\bottomrule
\end{tabular}
\kaynaknot{Taban olaydan önceki 12 ayın, \emph{Sonrası} sonraki üç ayın
ortalamasıdır; ikisi de 10.000 kelimede vukuat sayısıdır. \emph{Dönüş},
kaydın tabanın \%90'ına döndüğü ay sayısıdır. Korpus 31 Mart 2026'da bittiği
için son üç olayda 24 aylık pencere dolmuyor: oralarda dönüş gözlem dışıdır,
gerçekleşmediği anlamına gelmez.}
\end{table}


